\section{Introduction}

Knowledge-augmented predictors are usually built around a selection rule:
given a query and a pool of auxiliary contexts, the system decides which
contexts to pass to a fixed predictor. The rule is trained to rank candidates,
and deployment inherits the ranking as if it were a measurement. It is not. A
ranking model is optimised to order candidates correctly, not to report how
much a candidate will actually change the prediction. When the two disagree,
the router still selects the top-ranked candidate, and the prediction can end
up worse than if no context had been used at all.

The gap between ranking and measurement is not a small calibration detail. In
the deployment states we study, the standard deviation of the predicted
utility is roughly twice that of the true utility, and the absolute prediction
error is several times larger than the utility values it orders. Under those
conditions the top-ranked candidate is harmful in close to two of every five
decisions. A system that always routes therefore pays a hidden cost: it cannot
report how often it makes a prediction worse, and it has no mechanism for
declining to act.

We treat the decision to use a context as a decision that can be refused. The
router maintains a candidate score, converts it into a certified lower bound
through calibration on states it has already observed, and selects only when
the bound is positive. Every other decision becomes an abstention, including
the decision to stop adding contexts to a set that is already usable. This
reframes context selection from a ranking problem into a risk-controlled
decision problem, and it makes two quantities first class: the number of
harmful selections the router is willing to tolerate, and the fraction of
decisions on which it is still willing to act.

We formalise the resulting problem with an explicit harm budget. A harmful
selection is one whose true marginal utility is negative. The router
calibrates a threshold on a set of deployment states so that the expected
number of harmful selections per decision is at most a chosen level, and it
abstains whenever no candidate passes. Because the loss is monotone in the
threshold, conformal risk control supplies a finite-sample bound without
assumptions on the utility model or the data distribution, and the same
calibration yields the selective risk of the decisions that remain.

The natural alternative is interval certification: build a confidence interval
around each predicted utility and select only if the lower end is positive.
This is the rule that a distribution-free coverage guarantee suggests, and we
show that it collapses in exactly the regime that motivates the problem. When
the estimation error is larger than the utility itself, the interval radius
exceeds every candidate score, and the router abstains on almost every
decision; at the same harm budget, budget-controlled calibration retains
several times more decisions. Certification and control are therefore not
interchangeable, and the difference is measurable.

We call the resulting method Risk-Controlled Marginal-Utility Routing. It
reuses a response-aware utility model to score candidates, then adds the part
that decides whether to act: a conformal risk control threshold on deployment
states, an abstention rule, and a stopping rule for sequential selection. The
paper contributes the risk-controlled formulation of context selection, the
analysis of when interval certification fails and when a calibrated threshold
survives a change of dataset or of router training run, and an empirical study
that covers controlled harmful fractions, deployment states from six standard
traffic benchmarks, threshold transfer, and end-to-end selective routing.
